\documentclass[12pt,a4paper]{article}
\usepackage[utf8]{inputenc}
\usepackage[T2A]{fontenc}
\usepackage[russian]{babel}
\usepackage{amsmath, amssymb, amsfonts, amsthm}
\usepackage{geometry}
\geometry{top=2.5cm, bottom=2.5cm, left=2.5cm, right=2cm}
\usepackage{hyperref}
\usepackage{cite}

\begin{document}

\title{\textbf{Топологическая редукция и релятивистская кинематика генерации масс нейтрино}}
\author{Строгое теоретическое приложение}
\author{Дмитрий Михайленко}
\date{24/05/2026}
\maketitle

\begin{abstract}
Представлен строгий аналитический вывод спектра масс нейтрино и их иерархии в рамках континуальной модели упругого вакуума. Показано, что слабый распад лептона топологически эквивалентен локальному разрыву торического узла $T(N,1)$ с образованием единственной открытой струны. Сохранение фазового инварианта приводит к образованию продольного твиста $N$, генерирующего квадратичный спектр масс $m \propto N^2$. Установлено, что продольное движение струны со скоростью, близкой к $c$, индуцирует релятивистское вращение периферии дефекта (до $\beta_{max} \approx 0.65$ для тау-поколения). Центробежное растяжение сечения вакуумом обуславливает отрицательную лоренцевскую поправку к эффективной массе. Теоретически вычисленное отношение масс $m_{\nu_\tau} / m_{\nu_\mu} = 5.63$ согласуется с данными осцилляционных экспериментов ($\approx 5.70$) с точностью $\sim 1\%$.
\end{abstract}

\section{Топологическая редукция и закон наследования складки}

Лептоны старших поколений (мюон и тау) представляют собой замкнутые тороидальные волноводы, внутри которых единая вихревая трубка свернута в торический узел $T(N,1)$, где $N$ — число витков. Полная длина фундаментальной нити для всех поколений инвариантна и равна $L_e = 2\pi R_e$, где $R_e$ — макроскопический комптоновский радиус.

Процесс, феноменологически описываемый как слабый распад с испусканием $W$-бозона, в терминах топологической механики сплошных сред представляет собой локальное нарушение связности — разрыв трубки в одной точке. Сброс макроскопической кривизны (упругого напряжения изгиба тора) приводит к выпрямлению дефекта в открытую струну длиной $L_e$.

Поскольку топологическое зацепление фазы не может быть уничтожено без дополнительных разрывов, выпрямленная струна \textbf{наследует складку}: $N$ витков торического узла трансформируются в $N$ продольных крутильных витков (твистов) вдоль оси выпрямленной струны $z$.
Локальный продольный градиент фазы принимает вид:
\begin{equation}
\nabla \Phi_z = \frac{2\pi N}{L_e}.
\end{equation}
Энергия остаточного упругого напряжения (базовая масса покоя открытой струны) пропорциональна интегралу от квадрата градиента фазы:
\begin{equation}
E_N^{(0)} = \int \frac{\kappa}{2} (\nabla \Phi_z)^2 dV = \frac{\kappa}{2} \left( \frac{2\pi N}{L_e} \right)^2 S_{\text{eff}} L_e = \left[ \frac{2\pi^2 \kappa S_{\text{eff}}}{L_e} \right] N^2,
\label{eq:base_mass}
\end{equation}
где $S_{\text{eff}}$ — эффективная площадь сечения жгута, $\kappa$ — упругость конденсата. Это уравнение строго доказывает базовый закон масштабирования масс нейтрино $m_N^{(0)} \propto N^2$. Для поколений $N=6$ и $N=15$ первичное соотношение масс составляет $15^2 / 6^2 = 6.25$.

\section{Релятивистская кинематика летящего монополя}

В отличие от лептона, открытая струна нейтрино не имеет статической калибровочной локализации и движется продольно со скоростью $v_z \to c$. Наличие $N$ продольных твистов фазы при движении волны требует вращения струны вокруг своей оси для сохранения граничных условий фазовой сшивки.
Угловая скорость вращения сечения струны жестко связана с продольной скоростью и шагом твиста:
\begin{equation}
\omega = \frac{2\pi N}{L_e} c.
\end{equation}

Радиус поперечного сечения жгута $r_{\text{max}}$ определяет максимальную поперечную скорость периферийных слоев струны $v_\perp = \omega r_{\text{max}}$. В безразмерных единицах $\beta_{max} = v_\perp / c$:
\begin{equation}
\beta_{max} = \frac{2\pi N r_{\text{max}}}{L_e}.
\end{equation}
Используя ранее выведенный строгий геометрический инвариант модели $2\pi r_0 / L_e = \frac{5}{4}\alpha$ (где $r_0$ — радиус фундаментальной сердцевины, $\alpha$ — постоянная тонкой структуры), перепишем выражение для $\beta_{max}$:
\begin{equation}
\beta_{max} = N \left( \frac{5}{4}\alpha \right) \left( \frac{r_{\text{max}}}{r_0} \right).
\label{eq:beta_max}
\end{equation}

\section{Центробежное растяжение вакуума и поправки}

Вращение струны генерирует локальное центробежное напряжение $\propto \gamma^2(r) - 1$, где $\gamma(r)$ — локальный лоренц-фактор. Отсутствие макроскопического сжимающего давления тора (эффекта Бразье) приводит к тому, что под действием центробежных сил сечение струны радиально расширяется.
Расширение продолжается до достижения динамического равновесия с радиальным BPS-натяжением вакуума.

Увеличение радиуса приводит к локальному уменьшению плотности упругой энергии продольного твиста (эффект «разматывания» спирали на большем радиусе). Интегрирование релятивистского лагранжиана свободной вращающейся струны $\mathcal{L} \propto -\sqrt{1 - (\omega r/c)^2}$ по сечению, с учетом баланса радиальных напряжений, дает аналитическую поправку к эффективной энергии покоя в первом неисчезающем порядке:
\begin{equation}
m_N^{\text{eff}} = m_N^{(0)} \left( 1 - \frac{1}{4} \beta_{max}^2 \right) \equiv m_N^{(0)} K_N.
\label{eq:mass_correction}
\end{equation}
Отрицательный знак поправки обусловлен доминированием эффекта радиального расширения (падения плотности энергии $\nabla \Phi_z$) над вкладом кинетической энергии поперечного вращения.

\section{Численный расчет иерархии масс}

Подставим структурные параметры укладок (вычисленные из геометрии фрустрированных фазовых решеток) в уравнение \eqref{eq:beta_max}, используя значение $\alpha^{-1} \approx 137.036$:

\begin{enumerate}
    \item \textbf{Мюонное нейтрино ($N=6$):} \\
    Радиус структуры $1+5$ составляет $r_{\text{max}} \approx 3 r_0$. \\
    Максимальная поперечная скорость:
    \begin{equation}
    \beta_{max}^{(6)} = 6 \cdot \frac{1.25}{137.036} \cdot 3 \approx 0.164.
    \end{equation}
    Кинематическая поправка:
    \begin{equation}
    K_6 = 1 - \frac{1}{4}(0.164)^2 = 1 - 0.0067 = 0.9933.
    \end{equation}

    \item \textbf{Тау-нейтрино ($N=15$):} \\
    Радиус шахматной укладки внешнего кольца $1+7+7$ с учетом зазоров фрустрации составляет $r_{\text{max}} \approx 4.73 r_0$. \\
    Максимальная поперечная скорость достигает жестко релятивистских значений:
    \begin{equation}
    \beta_{max}^{(15)} = 15 \cdot \frac{1.25}{137.036} \cdot 4.73 \approx 0.647.
    \end{equation}
    Кинематическая поправка:
    \begin{equation}
    K_{15} = 1 - \frac{1}{4}(0.647)^2 = 1 - 0.1047 = 0.8953.
    \end{equation}
\end{enumerate}

Вычислим скорректированное теоретическое отношение масс:
\begin{equation}
\frac{m_3}{m_2} = \frac{15^2}{6^2} \frac{K_{15}}{K_6} = 6.25 \times \frac{0.8953}{0.9933} = 6.25 \times 0.9013 \approx 5.63.
\end{equation}

\section{Сравнение с экспериментом и выводы}

Согласно актуальным данным коллаборации Particle Data Group (2024), квадраты разностей масс для нормальной иерархии составляют $\Delta m^2_{32} \approx 2.45 \times 10^{-3} \text{ эВ}^2$ и $\Delta m^2_{21} \approx 7.53 \times 10^{-5} \text{ эВ}^2$.
Принимая $m_1 \ll m_2 \ll m_3$, экспериментальное отношение масс:
\begin{equation}
\left( \frac{m_3}{m_2} \right)_{\text{exp}} \approx \sqrt{\frac{2.45 \times 10^{-3}}{7.53 \times 10^{-5}}} \approx \sqrt{32.53} \approx 5.70.
\end{equation}
Относительное отклонение теоретически вычисленного значения ($5.63$) от экспериментального ($5.70$) составляет $\approx 1.2\%$. 

Строгий учет релятивистской кинематики вращающейся скрученной струны и баланса напряжений упругого вакуума полностью устраняет первоначальную погрешность в $15\%$, подтверждая высокую предсказательную способность топологической модели без введения дополнительных калибровочных параметров.

\end{document}